\section{Convolutional Sequence-to-Sequence Learning for Summarization}
\label{sec:topicconvabssum}


\begin{figure}[t!]
  \center{\footnotesize
    \begin{tikzpicture}[scale=0.45]
      
      \begin{scope}[shift={(0,0)}] 
        \draw [rounded corners=5pt] (-2.5,1.75) rectangle (15,-21.5);
        
        \draw [fill=sepia,opacity=0.1] (0,0) rectangle (1.5,-0.5);
        \draw [fill=sepia,opacity=0.3] (1.5,0) rectangle (6,-0.5);
        \draw [fill=sepia,opacity=0.1] (6,0) rectangle (7.5,-0.5);
        \draw [fill=orangered,opacity=0.1] (0,-0.5) rectangle (1.5,-1);
        \draw [fill=orangered,opacity=0.3] (1.5,-0.5) rectangle (6,-1);
        \draw [fill=orangered,opacity=0.1] (6,-0.5) rectangle (7.5,-1);
        
        \draw [gray] (0,-0.5) -- (7.5,-0.5);
        \draw [gray] (1.5,0) -- (1.5,-1);
        \draw [gray] (3,0) -- (3,-1);
        \draw [gray] (4.5,0) -- (4.5,-1);
        \draw [gray] (6,0) -- (6,-1);
        \node at (8.5,-0.75) {\tiny $x_i+p_i$};
        \node at (8.5,-0.25) {\tiny $t'_i \otimes t_D$};  
        \node at (-0.5,-0.5) {$\mathbf{e}$};

        \node at (1,0.75) {\tiny \rotatebox{45}{[PAD]}}; 
        \node at (2.5,0.75) {\tiny \rotatebox{45}{England}}; 
        \node at (4,0.75) {\tiny \rotatebox{45}{won}}; 
        \node at (5.25,0.75) {\tiny \rotatebox{45}{.}}; 
        \node at (7,0.75) {\tiny \rotatebox{45}{[PAD]}}; 
        
        % Convolution Part
        
        \node at (0,-2.5) {\tiny Convolutions};
        
        \draw [fill=midnightblue,opacity=0.2] (3.75,-1.5) -- (6.75, -1.5) -- (5.25, -3.5) -- (3.75,-1.5);
        \draw [fill=midnightblue,opacity=0.2] (2.25,-1.5) -- (5.25, -1.5) -- (3.75, -3.5) -- (2.25,-1.5);
        \draw [fill=midnightblue,opacity=0.2] (0.75,-1.5) -- (3.75, -1.5) -- (2.25, -3.5) -- (0.75,-1.5);
        
        \draw [fill=white,opacity=1] (1.75,-3.25) rectangle (2.75,-3.75);
        \draw [fill=white,opacity=1] (3.25,-3.25) rectangle (4.25,-3.75);
        \draw [fill=white,opacity=1] (4.75,-3.25) rectangle (5.75,-3.75);
        
        \draw (2.25,-3.25) -- (2.25,-3.75);
        \draw (3.75,-3.25) -- (3.75,-3.75);
        \draw (5.25,-3.25) -- (5.25,-3.75);
        
        % Gated Linear Unit

        \node at (0,-4.45) {\tiny GLU};
        
        \draw [->] (2,-3.75) -- (2,-4.25);
        \draw [gray] (2,-4.45) circle (0.2cm);
        \node at (2, -4.45) {\tiny $f$};
        \draw [->] (2.5,-3.75) -- (2.5,-5);
        \node at (2.5, -5.2) {$\otimes$};
        \draw [->] (2, -4.65) -- (2, -5.2) -- (2.25,-5.2);
        \draw [->] (2.5,-5.4) -- (2.5,-6);

        \draw [->] (3.5,-3.75) -- (3.5,-4.25);
        \draw [gray] (3.5, -4.45) circle (0.2cm);
        \node at (3.5, -4.45) {\tiny $f$};
        \draw [->] (4,-3.75) -- (4,-5);
        \node at (4,-5.2) {$\otimes$};
        \draw [->] (3.5,-4.65) -- (3.5,-5.2) -- (3.75,-5.2);
        \draw [->] (4,-5.4) -- (4,-6);

        \draw [->] (5,-3.75) -- (5,-4.25);
        \draw [gray] (5, -4.45) circle (0.2cm);
        \node at (5, -4.45) {\tiny $f$};
        \draw [->] (5.5,-3.75) -- (5.5,-5);
        \node at (5.5, -5.2) {$\otimes$};
        \draw [->] (5, -4.65) -- (5,-5.2) -- (5.25,-5.2);
        \draw [->] (5.5,-5.4) -- (5.5,-6);

        \draw [fill=midnightblue,opacity=0.2] (1.75,-6) rectangle (3.25,-7);
        \draw [fill=midnightblue,opacity=0.2] (3.25,-6) rectangle (4.75,-7);
        \draw [fill=midnightblue,opacity=0.2] (4.75,-6) rectangle (6.25,-7);
        \node at (1.25,-6.5) {$\mathbf{z^u}$};

        % Attention Part
        \node at (-1,-8.75) {\scriptsize \rotatebox{90}{Attention}}; 

        \draw [->] (2.5,-7) -- (2.5,-8);
        \draw [->] (4,-7) -- (4,-8);
        \draw [->] (5.5,-7) -- (5.5,-8);
        
        
        \draw [fill=forestgreen,opacity=0.2] (1.75,-8) rectangle (3.25,-8.5);
        \draw [fill=forestgreen,opacity=0.8] (3.25,-8) rectangle (4.75,-8.5);
        \draw [fill=forestgreen,opacity=0.7] (4.75,-8) rectangle (6.25,-8.5);
        \draw [fill=forestgreen,opacity=0.4] (1.75,-8.5) rectangle (3.25,-9);
        \draw [fill=forestgreen,opacity=0.8] (3.25,-8.5) rectangle (4.75,-9);
        \draw [fill=forestgreen,opacity=0.2] (4.75,-8.5) rectangle (6.25,-9);
        \draw [fill=forestgreen,opacity=0.1] (1.75,-9) rectangle (3.25,-9.5);
        \draw [fill=forestgreen,opacity=0.2] (3.25,-9) rectangle (4.75,-9.5);
        \draw [fill=forestgreen,opacity=0.3] (4.75,-9) rectangle (6.25,-9.5);

        % 
        \draw [fill=sepia,opacity=0.3] (10,-6) rectangle (14.5,-6.5);
        \draw [fill=orangered,opacity=0.3] (10,-6.5) rectangle (14.5,-7);
        \draw [gray] (11.5,-6) -- (11.5,-7);
        \draw [gray] (13,-6) -- (13,-7);
        \draw [fill=midnightblue,opacity=0.2] (10,-6) rectangle (14.5,-7);
        
        \draw [rounded corners=5pt,->] (9.5,-0.5) -- (12.25, -0.5) -- (12.25, -5.75);
        \draw [->] (6.25,-6.5) -- (9.4, -6.5);
        \node at (9.7,-6.5) {$\oplus$};
        % \node at (9.7,-6.5) {$\oplus$};
      \end{scope}

      \begin{scope}[shift={(-1.5,-19.5)}] 
        
        \draw [fill=sepia,opacity=0.1] (0,0) rectangle (4.5,0.5);
        \draw [fill=sepia,opacity=0.3] (4.5,0) rectangle (7.5,0.5);        
        \draw [fill=orangered,opacity=0.1] (0,0.5) rectangle (4.5,1);
        \draw [fill=orangered,opacity=0.3] (4.5,0.5) rectangle (7.5,1);

        \node at (0.75,-0.75) {\tiny \rotatebox{45}{[PAD]}}; 
        \node at (2.25,-0.75) {\tiny \rotatebox{45}{[PAD]}}; 
        \node at (3.75,-0.75) {\tiny \rotatebox{45}{[PAD]}}; 
        \node at (5.25,-0.75) {\tiny \rotatebox{45}{Match}}; 
        \node at (6.75,-0.75) {\tiny \rotatebox{45}{report}}; 
        

        \draw [gray] (0,0.5) -- (7.5,0.5);
        \draw [gray] (1.5,0) -- (1.5,1);
        \draw [gray] (3,0) -- (3,1);
        \draw [gray] (4.5,0) -- (4.5,1);
        \draw [gray] (6,0) -- (6,1);
        \node at (8.5,0.75) {\tiny $x'_i+p'_i$};
        \node at (8.5,0.25) {\tiny $t_D$};        
        \node at (-0.5,0.5) {$\mathbf{g}$};

        % Convolution Part
        
        \node at (1,3) {\tiny Convolutions};

        \draw [fill=midnightblue,opacity=0.3] (3.75,1.5) -- (6.75, 1.5) -- (6.75, 3.5) -- (3.75,1.5);
        \draw [fill=midnightblue,opacity=0.2] (2.25,1.5) -- (5.25, 1.5) -- (5.25, 3.5) -- (2.25,1.5);
        \draw [fill=midnightblue,opacity=0.1] (0.75,1.5) -- (3.75, 1.5) -- (3.75, 3.5) -- (0.75,1.5);
        
        \draw [fill=white,opacity=1] (3,3.25) rectangle (4,3.75);
        \draw [fill=white,opacity=1] (4.5,3.25) rectangle (5.5,3.75);
        \draw [fill=white,opacity=1] (6,3.25) rectangle (7,3.75);
        
        \draw (3.5,3.25) -- (3.5,3.75);
        \draw (5,3.25) -- (5,3.75);
        \draw (6.5,3.25) -- (6.5,3.75);
        
        % Gated Linear Unit
        \node at (1,4.45) {\tiny GLU};

        \draw [->] (3.25,3.75) -- (3.25,4.25);
        \draw [gray] (3.25, 4.45) circle (0.2cm);
        \node at (3.25, 4.45) {\tiny $f$};
        \draw [->] (3.75,3.75) -- (3.75,5);
        \node at (3.75, 5.2) {$\otimes$};
        \draw [->] (3.25, 4.65) -- (3.25, 5.2) -- (3.5,5.2);
        \draw [->] (3.75,5.4) -- (3.75,6);

        \draw [->] (4.75,3.75) -- (4.75,4.25);
        \draw [gray] (4.75, 4.45) circle (0.2cm);
        \node at (4.75, 4.45) {\tiny $f$};
        \draw [->] (5.25,3.75) -- (5.25,5);
        \node at (5.25, 5.2) {$\otimes$};
        \draw [->] (4.75, 4.65) -- (4.75,5.2) -- (5,5.2);
        \draw [->] (5.25,5.4) -- (5.25,6);

        \draw [->] (6.25,3.75) -- (6.25,4.25);
        \draw [gray] (6.25, 4.45) circle (0.2cm);
        \node at (6.25, 4.45) {\tiny $f$};
        \draw [->] (6.75,3.75) -- (6.75,5);
        \node at (6.75, 5.2) {$\otimes$};
        \draw [->] (6.25, 4.65) -- (6.25,5.2) -- (6.5,5.2);
        \draw [->] (6.75,5.4) -- (6.75,6);

        \draw [fill=midnightblue,opacity=0.1] (3,6) rectangle (4.5,7);
        \draw [fill=midnightblue,opacity=0.2] (4.5,6) rectangle (6,7);
        \draw [fill=midnightblue,opacity=0.3] (6,6) rectangle (7.5,7);
        
        \node at (2.5,6.5) {$\mathbf{h^{\ell}}$};

        % After Attention weights
        \draw [fill=forestgreen,opacity=0.3] (11.5,8) rectangle (13,9);
        \draw [fill=forestgreen,opacity=0.2] (13,8) rectangle (14.5,9);
        \draw [fill=forestgreen,opacity=0.3] (14.5,8) rectangle (16,9);     
        \draw [->] (12.25,8) -- (12.25,7);     
        \draw [->] (13.75,8) -- (13.75,7);     
        \draw [->] (15.25,8) -- (15.25,7); 
        \node at (11,8.5) {$\mathbf{c^{\ell}}$};

        \draw [fill=midnightblue!20!forestgreen,opacity=0.3] (11.5,6) rectangle (13,7);
        \draw [fill=midnightblue!60!forestgreen,opacity=0.2] (13,6) rectangle (14.5,7);
        \draw [fill=midnightblue!80!forestgreen,opacity=0.3] (14.5,6) rectangle (16,7);        
        \draw [->] (7.5,6.5) -- (10.9, 6.5);
        \node at (11.2,6.5) {$\oplus$}; 
        \node at (10.5,7) {$\mathbf{h^L}$};
        
        % Predict
        \draw [->] (12.25,6) -- (12.25,1); 
        \draw [fill=gray,opacity=0.3] (11.5,0) rectangle (13,1);
        \draw [->] (13.75,6) -- (13.75,1);     
        \draw [fill=gray,opacity=0.3] (13,0) rectangle (14.5,1);
        \draw [->] (15.25,6) -- (15.25,1);     
        \draw [fill=gray,opacity=0.3] (14.5,0) rectangle (16,1);
        
        \node at (12.25,-0.75) {\tiny \rotatebox{45}{Match}}; 
        \node at (13.75,-0.75) {\tiny \rotatebox{45}{report}}; 
        \node at (15.25,-0.75) {\tiny \rotatebox{45}{.}}; 
        

        % Connect to Attention
        \draw [rounded corners=5pt,->] (3.75,7) -- (3.75,7.5) -- (1.5, 7.5) -- (1.5, 11.25) -- (3.25, 11.25);
        \draw [rounded corners=5pt,->] (5.25,7) -- (5.25,8) -- (2, 8) -- (2, 10.75) -- (3.25, 10.75);
        \draw [rounded corners=5pt,->] (6.75,7) -- (6.75,8.5) -- (2.5, 8.5) -- (2.5, 10.25) -- (3.25, 10.25);
        % Predict using attention
        \draw [gray] (12.25,11.25) circle (0.2cm);
        \node at (12.25,11.25) {\fontsize{3}{5}\selectfont $\sum$};
        \draw [->] (7.75, 11.25) -- (12.05,11.25);
        \draw [->] (12.25, 11.05) -- (12.25,9);

        \draw [gray] (13.75,10.75) circle (0.2cm);
        \node at (13.75,10.75) {\fontsize{3}{5}\selectfont $\sum$};
        \draw [->] (7.75, 10.75) -- (13.5,10.75);
        \draw [->] (13.75, 10.5) -- (13.75,9);
        
        \draw [gray] (15.25,10.25) circle (0.2cm);
        \node at (15.25,10.25) {\fontsize{3}{5}\selectfont $\sum$};
        \draw [->] (7.75, 10.25) -- (15.05,10.25);
        \draw [->] (15.25, 10.05) -- (15.25,9);
        
        \node at (15,11.75) {$\Big\Downarrow$};
      \end{scope}
    \end{tikzpicture}
  }
  \vspace{-0.5cm}
  \caption{Topic-conditioned convolutional model for extreme
    summarization.}\label{fig:architecture}
\end{figure}

Unlike tasks like machine translation and paraphrase generation where
there is often a one-to-one semantic correspondence between source and
target words, document summarization must distill the content of the
document into a few important facts. This is even more challenging for
our task, where the compression ratio is extremely high, and pertinent
content can be easily missed.


%to distinguish content which is relevant for the
%summary from content which isn't. As such the challenges for extreme
%summarization are primarily in modeling long-range dependencies and
%recognizing pertinent content in the document.

Recently, a convolutional alternative to sequence modeling has been
proposed showing promise for machine translation
\cite{convenc_mt,convseq2seq} and story generation
\cite{fan-hier-gen}. We believe that convolutional architectures are
attractive for our summarization task for at least two
reasons. Firstly, contrary to recurrent networks which view the input
as a chain structure, convolutional networks can be stacked to
represent large context sizes. Secondly, hierarchical features can be
extracted over larger and larger contents, allowing to represent
long-range dependencies efficiently through shorter paths.

Our model builds on the work of \newcite{convseq2seq} who develop an
encoder-decoder architecture for machine translation with an attention
mechanism \cite{NIPS2015_5846} based exclusively on deep convolutional
networks. We adapt this model to our summarization task by allowing it
to recognize pertinent content (i.e.,~by foregrounding salient words
in the document).  In particular, we improve the convolutional encoder
by associating each word with a vector representing topic salience,
and the convolutional decoder by conditioning each word prediction on
the document topic vector.


% uses information about topical words in the document and the topic
% of the document to calibrate the modeling of long-range dependencies
% with the important content in the document for abstractive summary
% generation.

\paragraph{Model Overview}

At the core of our model is a simple convolutional block structure
that computes intermediate states based on a fixed number of input
elements. Our convolutional encoder (shown at the top of Figure
\ref{fig:architecture}) applies this unit across the document. We
repeat these operations in a stacked fashion to get a multi-layer
hierarchical representation over the input document where words at
closer distances interact at lower layers while distant words interact
at higher layers. The interaction between words through hierarchical
layers effectively captures long-range dependencies. 
% between distant
%words.% than when
%they are modelled with recurrent neural networks.

Analogously, our convolutional decoder (shown at the bottom of
Figure~\ref{fig:architecture}) uses the multi-layer convolutional
structure to build a hierarchical representation over what has been
predicted so far. Each layer on the decoder side determines useful
source context by attending to the encoder representation before it
passes its output to the next layer. This way the model remembers
which words it previously attended to and applies multi-hop attention
(shown at the middle of Figure \ref{fig:architecture}) per time
step. The output of the top layer is passed to a softmax classifier to
predict a distribution over the target vocabulary.


Our model assumes access to word and document topic
distributions. These can be obtained by any topic model, however we
use Latent Dirichlet Allocation (LDA; \citealt{Blei:2003:LDA}) in our
experiments; we pass the distributions obtained from LDA directly to
the network as additional input. This allows us to take advantage of
topic modeling without interfering with the computational advantages
of the convolutional architecture. The idea of capturing
document-level semantic information has been previously explored for
recurrent neural networks
\cite{mikolovZ12,Ghosh2016ContextualL,dieng-iclr17}, however, we are
not aware of any existing convolutional models.


\paragraph{Topic Sensitive Embeddings}

Let~$D$ denote a document consisting of a sequence of words $(w_1,
\ldots, w_m)$; we embed~$D$ into a distributional space $\mathbf{x} =
(x_1, \ldots, x_m)$ where $x_i \in \mathbb{R}^f$ is a column in
embedding matrix $M \in \mathbb{R}^{V\times f}$ (where $V$ is the
vocabulary size). We also embed the absolute word positions in the
document \mbox{$\mathbf{p} = (p_1, \ldots, p_m)$} where $p_i \in
\mathbb{R}^f$ is a column in position matrix $P \in
\mathbb{R}^{N\times f}$, and $N$ is the maximum number of
positions. Position embeddings have proved useful for convolutional
sequence modeling \cite{convseq2seq}, because, in contrast to RNNs,
they do not observe the temporal positions of words
\cite{D16-1248}. Let $t_D \in \mathbb{R}^{f'}$ be the topic
distribution of document~$D$ and $\mathbf{t'} = (t'_1, \ldots, t'_m)$
the topic distributions of words in the document (where $t'_i \in
\mathbb{R}^{f'}$). During encoding, we represent document~$D$ via
$\mathbf{e} = (e_1, \ldots, e_m)$, where~$e_i$ is:
\begin{align}
e_i = [(x_i+p_i);(t'_i\otimes t_D)] \in \mathbb{R}^{f+f'}, \label{eq:encoder}
\end{align}
\noindent and $\otimes$ denotes point-wise multiplication. The topic
distribution $t'_i$ of word $w_i$ essentially captures how topical the
word is in itself (local context), whereas the topic
distribution~$t_D$ represents the overall theme of the document
(global context). The encoder essentially enriches the context of the
word with its topical relevance to the document.


For every output prediction, the decoder estimates representation
$\mathbf{g} = (g_1, \ldots, g_n)$ for previously predicted words
$(w'_1, \ldots, w'_n)$ where $g_i$ is:
\begin{align}
g_i = [(x'_i+p'_i);t_D] \in \mathbb{R}^{f+f'}, \label{eq:decoder}
\end{align}
\noindent $x'_i$ and $p'_i$ are word and position embeddings of
previously predicted word~$w'_i$, and $t_D$ is the topic distribution
of the input document. Note that the decoder does not use the topic
distribution of~$w'_i$ as computing it on the fly would be
expensive. However, every word prediction is conditioned on the topic
of the document, enforcing the summary to have the same theme as the
document.

% We propose to use word topic information for local context and
% document topic to model global context.

\paragraph{Multi-layer Convolutional Structure}

Each convolution block, parametrized by $W \in \mathbb{R}^{2d\times
  kd}$ and $b_w \in \mathbb{R}^{2d}$, takes as input $X \in
\mathbb{R}^{k\times d}$ which is the concatenation of $k$ adjacent
elements embedded in a~$d$ dimensional space, applies one dimensional
convolution and returns an output element $Y \in \mathbb{R}^{2d}$. We
apply Gated Linear Units (GLU, $v:\mathbb{R}^{2d} \rightarrow
\mathbb{R}^{d}$, \citeauthor{pmlr-v70-dauphin17a}
\citeyear{pmlr-v70-dauphin17a}) on the output of the convolution
$Y$. Subsequent layers operate over the $k$ output elements of the
previous layer and are connected through residual connections
\cite{He2016DeepRL} to allow for deeper hierarchical
representation. We denote the output of the $\ell$th layer as
$\mathbf{h^{\ell}} = (h^{\ell}_1, \ldots, h^{\ell}_n)$ for the decoder
network, and $\mathbf{z^{\ell}} = (z^{\ell}_1, \ldots, z^{\ell}_m)$
for the encoder network.

\paragraph{Multi-hop Attention}

Our encoder and decoder are tied to each other through a multi-hop
attention mechanism. For each decoder layer $\ell$, we compute the
attention $a^{\ell}_{ij}$ of state $i$ and source element $j$ as:
\begin{align}
a^{\ell}_{ij} = \frac{\mbox{exp}(d^{\ell}_i \cdot z^u_j)}{\sum^m_{t=1} \mbox{exp}(d^{\ell}_i \cdot z^u_t)}, \label{eq:attention}
\end{align}
\noindent where $d^{\ell}_i = W^{\ell}_dh^{\ell}_i+b^{\ell}_i+g_i$ is
the decoder state summary combining the current decoder state
$h^{\ell}_i$ and the previous output element embedding $g_i$. The
vector $\mathbf{z^u}$ is the output from the last encoder layer
$u$. The conditional input $c^{\ell}_i$ to the current decoder layer
is a weighted sum of the encoder outputs as well as the input element
embeddings $e_j$:
\begin{align}
c^{\ell}_i = \sum^m_{j=1} a^{\ell}_{ij}(z^u_j+e_j). \label{eq:decinput}
\end{align}

The attention mechanism described here performs multiple attention
``hops'' per time step and considers which words have been previously
attended to. It is therefore different from single-step attention in
recurrent neural networks \cite{bahdanau-arxiv14}, where the attention
and weighted sum are computed over $\mathbf{z^u}$ only.

Our network uses multiple linear layers to project between the
embedding size $(f+f')$ and the convolution output size $2d$. They are
applied to~$\mathbf{e}$ before feeding it to the encoder, to the final
encoder output $\mathbf{z^u}$, to all decoder layers~$\mathbf{h^{\ell}}$
for the attention score computation, and to the final decoder output
$\mathbf{h^L}$ before the softmax.  We pad the input with~$k-1$ zero
vectors on both left and right side to ensure that the output of the
convolutional layers matches the input length. During decoding, we
ensure that the decoder does not have access to future information; we
start with $k$~zero vectors and shift the covolutional block to the
right after every prediction. The final decoder output $\mathbf{h^L}$
is used to compute the distribution over the target vocabulary $T$ as:
%
\begin{align}
  p(y_{i+1}|y_1,\ldots,y_i,D,t_D,\mathbf{t'}) = & \\
& \hspace*{-2cm}\mbox{softmax}(W_oh^L_i+b_o) \in \mathbb{R}^T. \label{eq:prediction}
\end{align}
%
We use layer normalization and weight initialization to
stabilize learning. 


Our topic-enhanced model calibrates long-range dependencies with
globally salient content. As a result, it provides a better
alternative to vanilla convolutional sequence models
\cite{convseq2seq} and RNN-based summarization models \cite{see-acl17}
for capturing cross-document inferences and paraphrasing.  At the same
time it retains the computational advantages of convolutional
models. Each convolution block operates over a fixed-size window of
the input sequence, allowing for simultaneous encoding of the input,
ease in learning due to the fixed number of non-linearities and
transformations for words in the input sequence.
%In
%contrast, RNNs are constrained by temporal dependencies in modeling
%long-range dependencies. Figure \ref{fig:architecture} gives an
%overview of our model.



% We refer the reader to \newcite{convseq2seq} for
% convolutional sequence modeling. Figure \ref{fig:architecture} gives
% an overview of our model.

%%% Local Variables:
%%% mode: latex
%%% TeX-master: "extreme-sum"
%%% End:

